%% ============================================================
%% Pré-Projeto de TCC — Agentic DAST
%% Ester Morais Neves — CEFET-MG — Março de 2026
%% Compilar com pdfLaTeX (Overleaf: compilador padrão)
%% ============================================================
\documentclass[
    12pt,
    openany,
    oneside,
    a4paper,
    brazil
]{abntex2}

%% ---- Codificação e idioma ------------------------------------
\usepackage[T1]{fontenc}
\usepackage[utf8]{inputenc}
\usepackage{lmodern}
\usepackage[brazil]{babel}

%% ---- Formatação geral ----------------------------------------
\usepackage{indentfirst}
\usepackage{microtype}
\usepackage{setspace}
\usepackage{geometry}

%% ---- Tabelas -------------------------------------------------
\usepackage{booktabs}
\usepackage{longtable}
\usepackage{array}
\usepackage{tabularx}
\usepackage{multirow}

%% ---- Matemática ----------------------------------------------
\usepackage{amsmath}
\usepackage{amssymb}

%% ---- Código-fonte --------------------------------------------
\usepackage{listings}
\usepackage{xcolor}

\definecolor{codeblue}{rgb}{0.0, 0.2, 0.6}
\definecolor{codegray}{rgb}{0.5,0.5,0.5}
\definecolor{codeorange}{rgb}{0.8, 0.4, 0.0}
\definecolor{backcolor}{rgb}{0.97,0.97,0.97}

\lstdefinestyle{tccstyle}{
    backgroundcolor=\color{backcolor},
    commentstyle=\color{codegray},
    keywordstyle=\bfseries\color{codeblue},
    numberstyle=\tiny\color{codegray},
    stringstyle=\color{codeorange},
    basicstyle=\small\ttfamily,
    breakatwhitespace=false,
    breaklines=true,
    captionpos=b,
    frame=single,
    keepspaces=true,
    numbers=left,
    numbersep=5pt,
    showspaces=false,
    showstringspaces=false,
    tabsize=2
}
\lstset{style=tccstyle, language=Java}

%% ---- Hyperlinks ----------------------------------------------
\usepackage{hyperref}
\hypersetup{
    pdftitle={Agentic DAST: Um Framework para Detecção de Vulnerabilidades em SPAs},
    pdfauthor={Ester Morais Neves},
    pdfsubject={Pré-projeto de Trabalho de Conclusão de Curso},
    pdfkeywords={DAST, SPA, LLM, segurança, vulnerabilidades},
    colorlinks=true,
    linkcolor=black,
    citecolor=black,
    urlcolor=blue,
    bookmarksdepth=4
}

%% ---- Referências (ABNT) --------------------------------------
\usepackage[brazilian,hyperpageref]{backref}
\usepackage[alf, abnt-emphasize=bf]{abntex2cite}

\renewcommand{\backrefpagesname}{Citado na(s) página(s):~}
\renewcommand{\backref}{}
\renewcommand*{\backrefalt}[4]{%
    \ifcase #1
        Nenhuma citação no texto.
    \or
        Citado na página #2.
    \else
        Citado #1 vezes nas páginas #2.
    \fi}

%% ============================================================
%% Metadados do documento
%% ============================================================
\titulo{Agentic DAST: Um Framework para Detecção Automatizada de
Vulnerabilidades em Single Page Applications
via Orquestração de LLMs e Navegação Semântica}

\autor{Ester Morais Neves}
\local{Belo Horizonte}
\data{Março de 2026}
\orientador{Prof. Mateus Felipe Tymburibá Ferreira}

\instituicao{%
  Centro Federal de Educação Tecnológica de Minas Gerais -- CEFET-MG
  \par
  Departamento de Engenharia de Computação
}

\tipotrabalho{Pré-projeto de Trabalho de Conclusão de Curso}

\preambulo{%
  Projeto de Conclusão de Curso apresentado ao Curso de Engenharia
  de Computação do Centro Federal de Educação Tecnológica de Minas
  Gerais, como requisito parcial à obtenção do título de Bacharel
  em Engenharia de Computação.
  \par
  Área de concentração: Segurança de Software
}

\makeindex

%% ============================================================
\begin{document}
%% ============================================================

\selectlanguage{brazil}
\frenchspacing

% ----- Elementos pré-textuais --------------------------------
\pretextual

\imprimircapa
\imprimirfolhaderosto

\pdfbookmark[0]{\contentsname}{toc}
\tableofcontents*
\cleardoublepage

% ----- Elementos textuais ------------------------------------
\textual

% =============================================================
\chapter{Introdução}
% =============================================================

As aplicações web contemporâneas migraram progressivamente do modelo
\textit{Multi-Page Application} (MPA), baseado em renderização no
servidor e navegação por hiperlinks estáticos, para o paradigma das
\textit{Single Page Applications} (SPAs). \textit{Frameworks} como
Angular, React e Vue.js promoveram uma inversão arquitetural na qual
toda a navegação é controlada por JavaScript no lado do cliente: o
estado da aplicação se altera sem recarregamento de página, e o
conteúdo é injetado dinamicamente no \textit{Document Object Model}
(DOM) por chamadas assíncronas a APIs (AJAX/\textit{Fetch})
\cite{stackhawk2024, getastra2024}.

Essa evolução, embora positiva para a experiência do usuário e para a
eficiência no desenvolvimento \textit{front-end}, criou uma lacuna
crítica nas metodologias de testes dinâmicos de segurança. Ferramentas
de \textit{Dynamic Application Security Testing} (DAST) consolidadas,
como o OWASP ZAP e \textit{scanners} comerciais legados, foram
projetadas para operar sobre o paradigma MPA, utilizando
\textit{crawlers} que analisam estaticamente o HTML em busca de
âncoras \texttt{<a>} para mapear a superfície de ataque
\cite{brightsec2024}. Diante de uma SPA moderna, esses
\textit{scanners} percebem uma única página opaca e são incapazes de
explorar os estados internos da aplicação --- exatamente onde residem
as vulnerabilidades mais críticas do ponto de vista de negócio.

As SPAs introduzem duas classes de desafios para ferramentas de DAST.
O primeiro é de \textbf{navegação}: sem renderização JavaScript, estados
profundos (carrinho, painel administrativo, fluxo de checkout) permanecem
inacessíveis. O segundo é de \textbf{interpretação semântica}: vulnerabilidades
como Lógica de Negócio (\textit{Business Logic Vulnerabilities} --- BLVs),
\textit{Cross-Site Scripting} (XSS) armazenado via campos dinâmicos,
ataques a \textit{JSON Web Tokens} (JWT) e exposição de dados sensíveis
em endpoints não documentados exigem compreensão do \textit{contexto} da
aplicação --- algo estruturalmente impossível para ferramentas baseadas em
correspondência de padrões sintáticos \cite{karadbhajne2025,portswigger2024}.
Exemplos canônicos de BLVs incluem: manipulação de parâmetros para comprar
itens com valor negativo, cupons de desconto em condições de corrida,
acesso não autorizado a recursos de outros usuários (IDOR) e escalada de
privilégios via adulteração de \textit{payload} JWT.

A emergência de Grandes Modelos de Linguagem (LLMs) com capacidades
avançadas de raciocínio --- como as séries Claude (Anthropic), GPT-4
(OpenAI) e Llama~3 (Meta) --- introduziu uma nova via para o problema.
LLMs são capazes de inferir a semântica de interfaces de usuário a
partir do conteúdo do DOM, raciocinar sobre o propósito de fluxos de
negócio e gerar hipóteses de ataque contextualmente relevantes. O
padrão ReAct (\textit{Reasoning + Acting}) \cite{yao2023react} e a
técnica de \textit{Chain-of-Thought} \cite{wei2022chain} demonstraram
que LLMs podem alternar entre raciocínio e execução de ações de forma
iterativa, comportamento diretamente aplicável à exploração autônoma de
aplicações web.

O presente trabalho propõe o \textbf{Agentic DAST}, um
\textit{framework} que integra o raciocínio semântico de LLMs com a
automação de navegadores via Puppeteer e o padrão Screenplay, com o
objetivo de superar estatisticamente as ferramentas DAST tradicionais
na detecção de vulnerabilidades em SPAs, cobrindo múltiplas categorias:
Lógica de Negócio (BLV), \textit{Broken Access Control}, XSS, JWT e
Exposição de Dados Sensíveis. A validação empírica será conduzida sobre
a aplicação OWASP Juice Shop \cite{owaspjuiceshop}, escolhida por combinar
arquitetura SPA realista (Angular + Node.js) com um catálogo documentado
de vulnerabilidades de múltiplas categorias, provendo \textit{Ground Truth}
confiável para o cálculo das métricas de eficácia.

% =============================================================
\chapter{Problema de Pesquisa}
% =============================================================

As ferramentas de DAST atuais falham sistematicamente na detecção de
vulnerabilidades de lógica de negócio em SPAs, pois carecem da
capacidade de navegar os estados internos dessas aplicações e de
compreender o contexto semântico dos fluxos de negócio. Em paralelo,
os avanços recentes em LLMs agênticos sugerem que tais sistemas podem
raciocinar sobre a intenção de interfaces de usuário e gerar hipóteses
de ataque contextualmente ricas --- mas a viabilidade e a eficácia
desse paradigma para DAST em SPAs não foram demonstradas empiricamente
na literatura.

Diante desse cenário, formula-se o seguinte \textbf{problema de
pesquisa}:

\begin{quote}
  \textit{Em que medida um framework de DAST baseado em agente autônomo
  LLM, orquestrando um navegador headless via padrão Screenplay,
  supera ferramentas tradicionais de DAST na detecção de vulnerabilidades
  em Single Page Applications, abrangendo categorias como Lógica de Negócio,
  Broken Access Control, XSS, JWT e Exposição de Dados Sensíveis?}
\end{quote}

% =============================================================
\chapter{Justificativa}
% =============================================================

\section{Dimensão Acadêmica}

A literatura científica sobre DAST para SPAs é escassa e concentra-se
predominantemente na caracterização das limitações de ferramentas
existentes, sem propor artefatos que as superem de forma empiricamente
mensurável \cite{brightsec2024, stackhawk2024}. Trabalhos recentes
sobre \textit{LLM-guided fuzzing} \cite{yang2025hybrid} demonstraram
resultados promissores para protocolos e linguagens de programação, mas
não abordaram o problema de navegar SPAs e detectar BLVs em contexto
web interativo. O Agentic DAST preenche essa lacuna com um artefato
tecnológico mensurável, validado por experimento controlado e
potencialmente publicável em conferências nacionais de Engenharia de
Software ou Segurança da Informação (SBES, SBSEG ou WebMedia).

\section{Dimensão Social e Profissional}

SPAs dominam o desenvolvimento web contemporâneo. A indústria
brasileira adota massivamente React e Angular em sistemas críticos:
portais bancários, plataformas de e-commerce, sistemas de saúde e
serviços governamentais. Uma ferramenta de DAST semanticamente
consciente tem aplicação direta em \textit{pipelines} de DevSecOps de
organizações de todos os portes, reduzindo o custo e o tempo de
detecção de vulnerabilidades que hoje dependem de auditorias manuais
dispendiosas. O \textit{framework} resultante será disponibilizado como
código aberto, democratizando o acesso a testes de segurança avançados.

\section{Dimensão Pessoal}

O interesse pela interseção entre segurança ofensiva e inteligência
artificial emergiu ao longo do curso de Engenharia de Computação,
especialmente após o estudo das disciplinas de Redes de Computadores e
Segurança da Informação. A possibilidade de aplicar raciocínio
computacional para automatizar uma tarefa hoje restrita a especialistas
humanos --- o \textit{pentest} de aplicações web --- representa um
desafio técnico e científico de alta relevância prática e pessoal.

% =============================================================
\chapter{Objetivos}
% =============================================================

\section{Objetivo Geral}

Desenvolver e validar experimentalmente o \textbf{Agentic DAST}, um
\textit{framework} que integra raciocínio semântico de LLMs com
automação de navegadores via Puppeteer e padrão Screenplay,
demonstrando superioridade estatística em métricas de \textit{Recall}
e F1-\textit{Score} na detecção de vulnerabilidades de lógica de
negócio em SPAs, em comparação com abordagens DAST tradicionais.

\section{Objetivos Específicos}

\begin{enumerate}
  \item \textbf{Realizar revisão sistemática} da literatura sobre:
  (i)~limitações de DAST em SPAs; (ii)~LLMs agênticos e padrão ReAct;
  (iii)~\textit{Chrome DevTools Protocol} (CDP) versus WebDriver;
  (iv)~padrão Screenplay; (v)~\textit{fuzzing} semântico guiado por
  LLMs.

  \item \textbf{Construir a Matriz de \textit{Ground Truth}} do OWASP
  Juice Shop, mapeando e classificando as vulnerabilidades de lógica de
  negócio do catálogo oficial como gabarito para o experimento
  comparativo.

  \item \textbf{Projetar e implementar os quatro módulos} do
  \textit{framework}: Módulo de Percepção (extrator DOM via CDP),
  Módulo Cognitivo (agente LLM com \textit{loop} ReAct e memória),
  Módulo Executor (Screenplay + Puppeteer com \textit{Abuse Cases}) e
  Módulo de Auditoria (OWASP ZAP \textit{daemon}).

  \item \textbf{Desenvolver biblioteca de \textit{Abuse Cases}} para as
  categorias de vulnerabilidades de lógica de negócio prioritárias
  (IDOR, manipulação de preço, condição de corrida, \textit{JWT
  forgery}, enumeração, \textit{information disclosure}), cobrindo no
  mínimo oito categorias distintas.

  \item \textbf{Implementar o \textit{Semantic Fuzzer}}, componente que
  utiliza LLM para gerar \textit{payloads} contextualmente relevantes
  a partir da análise semântica dos campos da interface, e avaliar sua
  superioridade em relação ao \textit{fuzzing} aleatório em cobertura
  de caminhos lógicos.

  \item \textbf{Conduzir experimento comparativo controlado} com três
  cenários --- (A) ZAP puro, (B) ZAP guiado por suíte E2E pré-gravada,
  (C) Agentic DAST autônomo em modo \textit{black-box} --- com três
  rodadas cada, analisando os resultados pelas métricas
  \textit{Recall}, \textit{Precision}, F1-\textit{Score} e Cobertura
  de \textit{Endpoints}.
\end{enumerate}

% =============================================================
\chapter{Questões de Pesquisa e Hipóteses}
% =============================================================

\section{Questões de Pesquisa}

\begin{description}
  \item[\textbf{QP1}] O Agentic DAST autônomo (Cenário~C) obtém
  \textit{Recall} significativamente superior ao OWASP ZAP tradicional
  (Cenário~A) na detecção de vulnerabilidades de lógica de negócio no
  OWASP Juice Shop?

  \item[\textbf{QP2}] A capacidade de navegação semântica do Agentic
  DAST resulta em maior cobertura de \textit{endpoints} e estados de UI
  do OWASP Juice Shop em comparação com os \textit{baselines}?

  \item[\textbf{QP3}] O \textit{Semantic Fuzzer} baseado em LLM gera
  \textit{payloads} que cobrem caminhos lógicos distintos dos cobertos
  pelo \textit{fuzzing} aleatório, produzindo detecções adicionais?
\end{description}

\section{Hipóteses Testáveis}

\begin{description}
  \item[\textbf{H1}] O Agentic DAST (Cenário~C) detecta, em média, um
  número estatisticamente significativamente maior de vulnerabilidades
  de lógica de negócio verdadeiro-positivas do que o OWASP ZAP puro
  (Cenário~A), medido pelo \textit{Recall}.

  \textit{Critério de confirmação:} $\text{Recall}(C) > \text{Recall}(A)$
  com diferença estatisticamente significativa ($p < 0{,}05$, teste de
  Wilcoxon para amostras emparelhadas, $n = 3$ rodadas).

  \item[\textbf{H2}] A cobertura de \textit{endpoints} da aplicação
  alvo pelo Agentic DAST é superior à dos Cenários~A e~B, demonstrando
  que a navegação guiada por LLM explora mais estados da SPA.

  \textit{Critério de confirmação:}
  $\text{Endpoints}(C) > \text{Endpoints}(A)$ e
  $\text{Endpoints}(C) > \text{Endpoints}(B)$.

  \item[\textbf{H3}] Os \textit{payloads} gerados pelo \textit{Semantic
  Fuzzer} ativam, em média, mais vulnerabilidades de lógica de negócio
  por \textit{payload} testado do que \textit{payloads} gerados por
  \textit{fuzzing} aleatório sobre os mesmos campos.

  \textit{Critério de confirmação:} Taxa de detecção semântica $>$
  taxa de detecção aleatória, avaliada quantitativamente e
  qualitativamente sobre amostras pareadas de \textit{payloads}.
\end{description}

% =============================================================
\chapter{Fundamentação Teórica Preliminar}
% =============================================================

Esta seção apresenta os pilares conceituais que embasam o
\textit{framework} proposto. Uma revisão sistemática aprofundada será
conduzida na Fase~1 do projeto.

\section{DAST em SPAs e Vulnerabilidades de Lógica de Negócio}

O \textit{Dynamic Application Security Testing} (DAST) testa
aplicações em tempo de execução, sem acesso ao código-fonte. As
ferramentas tradicionais baseiam-se em \textit{spiders} que percorrem
hiperlinks HTML estáticos, tornando-se cegas diante do roteamento
\textit{client-side} das SPAs \cite{stackhawk2024}. Vulnerabilidades de
Lógica de Negócio diferem das vulnerabilidades técnicas clássicas
(injeção SQL, XSS) por exigirem compreensão do domínio da aplicação
para serem detectadas \cite{portswigger2024, karadbhajne2025}.

\section{Chrome DevTools Protocol e Automação de Navegadores}

O \textit{Chrome DevTools Protocol} (CDP) provê uma interface
bidirecional de baixo nível com o motor Chromium/Blink via WebSockets,
permitindo intercepção de rede em tempo real, instrumentação de
JavaScript e isolamento de sessão com granularidade inacessível ao
Selenium WebDriver \cite{lightpanda2024}. O Puppeteer expõe o CDP de
forma quase direta, sendo o instrumento ideal para testes de segurança
ofensiva que exigem manipulação de requisições HTTP em nível de
protocolo.

\section{Padrão Screenplay}

O Screenplay Pattern \cite{serenityscreenplay} modela a automação de
testes em termos de cinco elementos fundamentais: Atores, Habilidades,
Tarefas, Interações e Perguntas. Ao desacoplar a intenção (Tarefa) da
implementação (Interação), permite que o Módulo Cognitivo raciocine em
termos semânticos de alto nível, sem necessidade de conhecer os
seletores CSS exatos da aplicação, que são resolvidos em tempo de
execução pelo Módulo Executor.

\section{LLMs Agênticos e o Padrão ReAct}

\citeonline{yao2023react} introduziram o padrão ReAct
(\textit{Reasoning + Acting}), no qual o LLM intercala passos de
raciocínio (\textit{thought}), ações executáveis (\textit{action}) e
observações do ambiente (\textit{observation}), produzindo trajetórias
de decisão interpretáveis. \citeonline{wei2022chain} demonstraram que
\textit{Chain-of-Thought prompting} habilita LLMs a resolver problemas
de múltiplas etapas com eficácia superior a abordagens diretas. Esses
paradigmas formam a base algorítmica do Módulo Cognitivo do Agentic
DAST.

\section{Fuzzing Semântico Guiado por LLMs}

\textit{Fuzzing} semântico avança o \textit{fuzzing} tradicional ao
gerar entradas que respeitam a gramática e o contexto da aplicação, mas
violam suas regras lógicas. \citeonline{yang2025hybrid} demonstraram
que LLMs com \textit{constraint solving} podem gerar \textit{payloads}
que cobrem caminhos de execução profundos inacessíveis a
\textit{fuzzers} por mutação aleatória. O \textit{Semantic Fuzzer} do
Agentic DAST aplica esse princípio ao contexto de campos de formulário
em SPAs.

% =============================================================
\chapter{Metodologia}
% =============================================================

A pesquisa adota a abordagem de \textbf{Design Science Research} (DSR)
\cite{hevner2004design}, que combina o desenvolvimento de um artefato
tecnológico com a geração de conhecimento científico rigoroso sobre sua
eficácia. O trabalho é dividido entre TCC~I e TCC~II, em sete fases
sequenciais, apresentadas na Tabela~\ref{tab:fases}.

% ----- Tabela de fases ----------------------------------------
\begin{longtable}{p{0.6cm}p{1.6cm}p{2.0cm}p{5.2cm}p{3.0cm}}
\caption{Fases do projeto e entregáveis.\label{tab:fases}}\\
\toprule
\textbf{F} & \textbf{Sem.} & \textbf{Período} &
\textbf{Atividades Principais} & \textbf{Entregável} \\
\midrule
\endfirsthead
\multicolumn{5}{l}{\textit{(Continuação da Tabela~\ref{tab:fases})}}\\
\toprule
\textbf{F} & \textbf{Sem.} & \textbf{Período} &
\textbf{Atividades Principais} & \textbf{Entregável} \\
\midrule
\endhead
\midrule
\multicolumn{5}{r}{\textit{Continua na próxima página}}\\
\endfoot
\bottomrule
\endlastfoot

F1 & S01–S07 & Mar–Abr 2026 &
Revisão sistemática (DAST em SPAs, LLMs agênticos, CDP, Screenplay,
fuzzing semântico); Construção da Matriz de Ground Truth &
Relatório Estado da Arte + Matriz GT \\[4pt]

F2 & S08–S11 & Abr–Mai 2026 &
Provisionamento Docker Compose: Juice Shop, ZAP daemon, Node.js
+ Puppeteer + Serenity/JS; estruturação do repositório &
Infraestrutura operacional \\[4pt]

F3 & S12–S19 & Mai–Jul 2026 &
Módulo Executor: Atores, Habilidades, Tarefas, primeiros Abuse Cases
(IDOR, preço negativo, feedback inválido) &
Biblioteca de Atores + 3 Abuse Cases \\[4pt]

F4 & S20–S27 & Jul–Set 2026 &
Módulo Cognitivo: loop ReAct, memória de curto/longo prazo, prompt
engineering iterativo, self-healing; integração Claude Haiku 4.5 /
Sonnet 4.6 &
Agente autônomo funcional \\[4pt]

F5 & S28–S31 & Set–Out 2026 &
Abuse Cases completos (JWT forgery, condição de corrida, enumeração);
Semantic Fuzzer &
8+ Abuse Cases + Semantic Fuzzer \\[4pt]

F6 & S32–S35 & Out 2026 &
Experimento comparativo (Cenários A, B, C); coleta e análise
estatística &
Tabelas VP/FP/FN + gráficos P/R/F1 \\[4pt]

F7 & S36 & Nov 2026 &
Redação do TCC (normas ABNT); submissão à banca; adaptação para artigo
científico &
TCC entregue + artigo submetido \\

\end{longtable}

\section{Protocolo Experimental}

O experimento comparativo (Fase~6) seguirá o protocolo a seguir.

\begin{itemize}
  \item \textbf{Objeto de estudo:} OWASP Juice Shop, versão estável
  mais recente, instância Docker isolada.

  \item \textbf{Cenário A --- Baseline Tradicional:} OWASP ZAP~2.x com
  \textit{spider} padrão, sem configuração adicional.

  \item \textbf{Cenário B --- Baseline E2E:} OWASP ZAP com
  \textit{spider} + suíte de testes funcionais E2E (Serenity/JS,
  Caminho Feliz) pré-gravada guiando o \textit{proxy}.

  \item \textbf{Cenário C --- Experimental:} Agentic DAST em modo
  \textit{black-box}, sem conhecimento prévio da estrutura da
  aplicação.

  \item \textbf{Repetições:} Três rodadas por cenário ($n = 3$), com
  reinicialização do banco de dados do Juice Shop entre rodadas para
  garantir estado inicial idêntico.

  \item \textbf{Análise estatística:} Média e desvio-padrão por
  cenário; teste de Wilcoxon (amostras emparelhadas) para comparação
  entre cenários, com $\alpha = 0{,}05$.
\end{itemize}

\section{Métricas de Avaliação}

A Tabela~\ref{tab:metricas} apresenta as métricas a serem coletadas e
calculadas para cada cenário.

\begin{table}[htb]
\centering
\caption{Métricas de avaliação do experimento comparativo.%
\label{tab:metricas}}
\begin{tabularx}{\textwidth}{lX}
\toprule
\textbf{Métrica} & \textbf{Descrição} \\
\midrule
Verdadeiro Positivo (VP) &
  Vulnerabilidade presente na Matriz GT detectada corretamente \\[2pt]
Falso Positivo (FP) &
  Alerta gerado para algo que não constitui vulnerabilidade \\[2pt]
Falso Negativo (FN) &
  Vulnerabilidade presente na Matriz GT não detectada \\[2pt]
\textit{Recall} &
  $\text{VP} / (\text{VP} + \text{FN})$ --- taxa de cobertura das BLVs
  conhecidas \\[2pt]
\textit{Precision} &
  $\text{VP} / (\text{VP} + \text{FP})$ --- taxa de alertas verdadeiros
  \\[2pt]
F1-\textit{Score} &
  $2 \cdot (\text{Precision} \cdot \text{Recall}) /
  (\text{Precision} + \text{Recall})$ \\[2pt]
Cobertura de \textit{Endpoints} &
  Número de rotas de API e estados de UI descobertos \\
\bottomrule
\end{tabularx}
\end{table}

\section{Ameaças à Validade}

\textbf{Validade interna:} A natureza probabilística dos LLMs pode
introduzir variação entre rodadas. Será controlada via temperatura fixa
($T = 0{,}3$), mesma versão de modelo e API, e análise do
desvio-padrão entre as três rodadas.

\textbf{Validade externa:} Os resultados são obtidos sobre o OWASP
Juice Shop (aplicação propositalmente vulnerável); a generalização para
sistemas reais exige validação adicional em ambientes de
\textit{staging} corporativo, prevista como entregável da Fase~7.

\textbf{Validade de construção:} A Matriz de \textit{Ground Truth} que
define o que é e o que não é uma vulnerabilidade de lógica de negócio
no Juice Shop será construída a partir do catálogo oficial da OWASP,
reduzindo viés de seleção subjetivo.

% =============================================================
\chapter{Infraestrutura Necessária}
% =============================================================

\section{Hardware e Software}

\begin{itemize}
  \item Computador pessoal com mínimo de 16~GB de RAM e 30~GB livres
  para containers Docker.
  \item Docker Desktop para orquestração dos três containers: Juice
  Shop, ZAP e Node.js.
  \item Node.js 20 LTS e Git para desenvolvimento e versionamento.
  \item VS Code com GitHub Copilot, disponível via GitHub Student
  Developer Pack.
\end{itemize}

\section{APIs de LLM}

\begin{itemize}
  \item \textbf{Claude API (Anthropic):} Haiku 4.5 (classificação
  rápida) e Sonnet 4.6 (raciocínio principal) com
  \textit{Prompt Caching} automático (redução de até 90\% no custo de
  \textit{tokens} repetidos). LLM principal do projeto
  \cite{anthropicdocs}.

  \item \textbf{Google AI Studio:} Gemini 2.5 Flash-Lite (1.000
  req/dia gratuitas) como \textit{fallback} sem custo.

  \item \textbf{Groq Cloud:} Llama 3.1 8B Instant (14.400 req/dia
  gratuitas) para iterações de alta frequência.

  \item \textbf{Ollama (local):} Llama 3 8B para desenvolvimento e
  depuração 100\% \textit{offline}, sem consumo de cotas.
\end{itemize}

\section{Dependências e Limitações Reconhecidas}

O trabalho não depende de participantes externos: todos os experimentos
são executados em ambiente Docker isolado e controlado. Requer
orientação periódica do orientador e aprovação do protocolo
experimental antes da execução dos cenários comparativos (Fase~6). A
qualidade dos resultados depende parcialmente dos \textit{prompts} de
sistema e da versão do LLM utilizado; variações nas cotas gratuitas das
APIs podem exigir ajustes no roteador de modelos durante os
experimentos.

% =============================================================
\chapter{Resultados Esperados}
% =============================================================

Ao término do trabalho, espera-se obter:

\begin{enumerate}
  \item \textbf{\textit{Framework} funcional e de código aberto}
  (Agentic DAST), disponibilizado em repositório Git público com
  instruções de instalação, configuração e exemplos de uso.

  \item \textbf{Evidência empírica de superioridade} do Agentic DAST
  em \textit{Recall} e F1-\textit{Score} para BLVs, validada por
  experimento controlado (3 execuções por cenário, análise estatística).

  \item \textbf{Biblioteca de \textit{Abuse Cases} reutilizável} para
  SPAs, modelada com Screenplay Pattern, cobrindo no mínimo oito
  categorias de vulnerabilidades de lógica de negócio.

  \item \textbf{\textit{Semantic Fuzzer}} demonstravelmente superior ao
  \textit{fuzzing} aleatório na cobertura de caminhos lógicos, com
  análise qualitativa das cadeias de raciocínio do LLM.

  \item \textbf{Artigo científico} submetido a conferência nacional de
  Engenharia de Software ou Segurança da Informação (SBES, SBSEG ou
  WebMedia).

  \item \textbf{Relatório de aplicação} do \textit{framework} em
  ambiente de \textit{staging} corporativo, demonstrando
  transferibilidade para contextos industriais.
\end{enumerate}

% =============================================================
\chapter{Cronograma}
% =============================================================

O projeto será desenvolvido ao longo de \textbf{36 semanas},
distribuídas entre TCC~I (Fases 1--3, março a julho de 2026) e TCC~II
(Fases 4--7, julho a novembro de 2026). A Tabela~\ref{tab:cronograma}
apresenta os marcos principais do projeto.

\begin{table}[htb]
\centering
\caption{Marcos e datas-chave do projeto.\label{tab:cronograma}}
\begin{tabular}{lcc}
\toprule
\textbf{Marco} & \textbf{Data} & \textbf{Semana} \\
\midrule
Início do projeto                              & 09/03/2026  & S01      \\
Infraestrutura Docker pronta                   & 27/04/2026  & S08      \\
Entrega TCC~I (Módulo Executor + Ground Truth) & Jul. 2026   & S19      \\
Agente autônomo funcionando (demo)             & 07/09/2026  & S27      \\
Experimentos concluídos                        & 05/10/2026  & S31      \\
Início da redação formal                       & 06/10/2026  & S32      \\
Entrega final TCC~II                           & 16/11/2026  & S36      \\
Submissão do artigo científico                 & 30/11/2026  & Pós-S36  \\
\bottomrule
\end{tabular}
\end{table}

% ----- Pós-textuais ------------------------------------------
\postextual

\bibliography{referencias}

\end{document}
